\begin{trapbox}
	\begin{description}[style=nextline, leftmargin=2.8em, labelsep=0.5em, itemsep=0.35em]
		\item[\textbf{\textcolor{CTrap}{易错点一（忽略零项隐患）}}] 直接假设 $a_n\neq 0$ 而不验证，若存在 $k$ 使 $a_k=0$，则倒数变换 $b_k$ 无定义。
		\item[\textbf{\textcolor{CTrap}{正确理解（先验非零条件）：}}] 使用本结论前必须确认对一切 $n$ 都有 $a_n\neq 0$，或通过递推初始项判断；若出现零项应分段处理。
		\item[\textbf{\textcolor{CTrap}{错因分析（缺少前置检查）：}}] 学生习惯直接套公式，忽略了倒数变换的前提，导致解题步骤中出现无意义表达式。
	\end{description}

	\medskip
	\begin{description}[style=nextline, leftmargin=2.8em, labelsep=0.5em, itemsep=0.35em]
		\item[\textbf{\textcolor{CTrap}{易错点二（错用递推类型）}}] 遇到 $a_{n+1}=\dfrac{p a_n+q}{r a_n + s}$（$q\neq 0$）仍然直接取倒数，误以为能化为线性。
		\item[\textbf{\textcolor{CTrap}{正确理解（区分分子结构）：}}] 本结论要求分子只有 $a_n$ 的一次项且无常数项；若分子有常数项 $q$，需先通过平移或取倒数后再做其他处理，不能简单套用。
		\item[\textbf{\textcolor{CTrap}{错因分析（形式类比错误）：}}] 学生看到分式形似就盲目使用，未检查分子是否满足纯一次，导致转换后不是线性递推，无法求解。
	\end{description}

	\medskip
	\begin{description}[style=nextline, leftmargin=2.8em, labelsep=0.5em, itemsep=0.35em]
		\item[\textbf{\textcolor{CTrap}{易错点三（分式化简失误）}}] 将 $a_n=\frac1{b_n}$ 代入后，分母通分或约简时出错，得到错误的 $b_{n+1}$ 表达。
		\item[\textbf{\textcolor{CTrap}{正确理解（规范通分步骤）：}}] 严格按 $a_{n+1}=\frac{f\cdot\frac1{b_n}}{g\cdot\frac1{b_n}+h}$ 化简为 $\frac{f}{g+hb_n}$ 的步骤进行，不跳步。
		\item[\textbf{\textcolor{CTrap}{错因分析（代数基本功弱）：}}] 分式运算复杂时容易出错，特别是忘记翻转相乘或约分时忽略分子分母整体性。
	\end{description}
\end{trapbox}
